% GENERATED FILE. Edit canonical lessons, then run npm run generate:books.
% canonical-source-hash: fnv1a64:be88b2b18abc7594
% canonical-lessons: HI-C482-qad, HI-C482-huliya, HI-C482-savla, HI-C482-gora, HI-C482-dubla

\chapter{How Someone Looks}
\label{ch:hi-how-someone-looks}
\hlchaptermodality{482}

\begin{chapteropening}
\textbf{By the end of this chapter:} \emph{I can describe how tall or thin someone is.}

Complete the last lesson of chapter 482: I can describe how tall or thin someone is.
\end{chapteropening}

\section[\texorpdfstring{\hi{क़द} (qad)}{क़द (qad)}]{\hi{क़द} (qad) — height; stature}

\label{lesson:HI-C482-qad}



\begin{quote}
Before the new one: say the Hindi for to stand on tiptoe, then the Hindi for to jump.
\end{quote}

\subsection*{You'll want to know: \hi{क़द}}
\textbf{\hi{क़द}} — \emph{qad} — \textquotedblleft{}height; stature\textquotedblright{}.

\begin{quote}
\hi{अरुण} \hi{का} \hi{क़द} \hi{अपने} \hi{पिता} \hi{से} \hi{लंबा} \hi{है।}
\end{quote}

\begin{quote}
\emph{Aruṇ kā qad apne pitā se lambā hai.}
\end{quote}

\begin{quote}
Arun is taller than his father.
\end{quote}

Masculine.

\begin{cousinweb}
\textbf{\hi{क़द}} (\emph{qad}) is usually traced to Arabic \emph{qadd}, \textquotedblleft{}stature, height\textquotedblright{}, which reached Hindi through Persian.
\end{cousinweb}

\begin{grammarlens}[title={What you've built}]
One more word for describing what you lived and felt.
\end{grammarlens}

\subsection*{Your turn}
\begin{itemize}
\raggedright
  \item \emph{Say it:} \emph{qad}
  \item \emph{Say it:} \emph{qad}, once more
  \item \emph{Say it:} \emph{Aruṇ kā qad apne pitā se lambā hai.} — Arun is taller than his father.
  \item \emph{Recall:} say the Hindi for to stand on tiptoe, then the Hindi for to jump, then say \emph{qad} again
\end{itemize}

\begin{tcolorbox}[breakable,colback=teal!4,colframe=teal!35!black,title={Before you move on}]
What does \hi{क़द} mean? (\textquotedblleft{}Height; stature\textquotedblright{}.) Say it once more.
\end{tcolorbox}
\section[\texorpdfstring{\hi{हुलिया} (huliyā)}{हुलिया (huliyā)}]{\hi{हुलिया} (huliyā) — appearance; looks}

\label{lesson:HI-C482-huliya}



\begin{quote}
Before the new one: say the Hindi for to jump, then the Hindi for height.
\end{quote}

\subsection*{You'll want to know: \hi{हुलिया}}
\textbf{\hi{हुलिया}} — \emph{huliyā} — \textquotedblleft{}appearance; looks\textquotedblright{}.

\begin{quote}
\hi{दस} \hi{साल} \hi{बाद} \hi{दोस्त} \hi{का} \hi{हुलिया} \hi{बिलकुल} \hi{बदल} \hi{गया} \hi{था।}
\end{quote}

\begin{quote}
\emph{das sāl bād dost kā huliyā bilkul badal gayā thā.}
\end{quote}

\begin{quote}
Ten years on, my friend's looks had changed completely.
\end{quote}

Masculine.

\begin{cousinweb}
\textbf{\hi{हुलिया}} (\emph{huliyā}) is usually traced to Arabic \emph{ḥulya}, \textquotedblleft{}ornament\textquotedblright{}, and so \textquotedblleft{}a description of someone's features\textquotedblright{}, which reached Hindi through Persian.
\end{cousinweb}

\begin{grammarlens}[title={What you've built}]
One more word for describing what you lived and felt.
\end{grammarlens}

\subsection*{Your turn}
\begin{itemize}
\raggedright
  \item \emph{Say it:} \emph{huliyā}
  \item \emph{Say it:} \emph{huliyā}, once more
  \item \emph{Say it:} \emph{das sāl bād dost kā huliyā bilkul badal gayā thā.} — Ten years on, my friend's looks had changed completely.
  \item \emph{Recall:} say the Hindi for to jump, then the Hindi for height, then say \emph{huliyā} again
\end{itemize}

\begin{tcolorbox}[breakable,colback=teal!4,colframe=teal!35!black,title={Before you move on}]
What does \hi{हुलिया} mean? (\textquotedblleft{}Appearance; looks\textquotedblright{}.) Say it once more.
\end{tcolorbox}
\section[\texorpdfstring{\hi{साँवला} (sā̃vlā)}{साँवला (sā̃vlā)}]{\hi{साँवला} (sā̃vlā) — dusky; dark-skinned}

\label{lesson:HI-C482-savla}



\begin{quote}
Before the new one: say the Hindi for height, then the Hindi for appearance.
\end{quote}

\subsection*{You'll want to know: \hi{साँवला}}
\textbf{\hi{साँवला}} — \emph{sā̃vlā} — \textquotedblleft{}dusky; dark-skinned\textquotedblright{}.

\begin{quote}
\hi{मेरी} \hi{बहन} \hi{का} \hi{रंग} \hi{साँवला} \hi{है।}
\end{quote}

\begin{quote}
\emph{merī bahan kā raṅg sā̃vlā hai.}
\end{quote}

\begin{quote}
My sister has a dusky complexion.
\end{quote}

\begin{cousinweb}
\textbf{\hi{साँवला}} (\emph{sā̃vlā}) is a tadbhava: CDIAL lists Hindi \emph{sā̃wlā}, \textquotedblleft{}brown, dark-complexioned\textquotedblright{}, under Sanskrit \emph{śyāmalá}, \textquotedblleft{}dark\textquotedblright{} (CDIAL 12665). Like \emph{gorā}, it describes a complexion and nothing more.
\end{cousinweb}

\begin{grammarlens}[title={What you've built}]
One more word for describing what you lived and felt.
\end{grammarlens}

\subsection*{Your turn}
\begin{itemize}
\raggedright
  \item \emph{Say it:} \emph{sā̃vlā}
  \item \emph{Say it:} \emph{sā̃vlā}, once more
  \item \emph{Say it:} \emph{merī bahan kā raṅg sā̃vlā hai.} — My sister has a dusky complexion.
  \item \emph{Recall:} say the Hindi for height, then the Hindi for appearance, then say \emph{sā̃vlā} again
\end{itemize}

\begin{tcolorbox}[breakable,colback=teal!4,colframe=teal!35!black,title={Before you move on}]
What does \hi{साँवला} mean? (\textquotedblleft{}Dusky; dark-skinned\textquotedblright{}.) Say it once more.
\end{tcolorbox}
\section[\texorpdfstring{\hi{गोरा} (gorā)}{गोरा (gorā)}]{\hi{गोरा} (gorā) — fair-skinned}

\label{lesson:HI-C482-gora}



\begin{quote}
Before the new one: say the Hindi for appearance, then the Hindi for dusky.
\end{quote}

\subsection*{You'll want to know: \hi{गोरा}}
\textbf{\hi{गोरा}} — \emph{gorā} — \textquotedblleft{}fair-skinned\textquotedblright{}.

\begin{quote}
\hi{मीरा} \hi{का} \hi{भाई} \hi{गोरा} \hi{है} \hi{और} \hi{मीरा} \hi{साँवली} \hi{है।}
\end{quote}

\begin{quote}
\emph{Mīrā kā bhāī gorā hai aur Mīrā sā̃vlī hai.}
\end{quote}

\begin{quote}
Mira's brother is fair and Mira is dusky.
\end{quote}

\begin{cousinweb}
\textbf{\hi{गोरा}} (\emph{gorā}) is a tadbhava: CDIAL lists Hindi \emph{gorā}, \textquotedblleft{}fair\textquotedblright{}, under Sanskrit \emph{gaurá}, \textquotedblleft{}white, yellowish, pale red\textquotedblright{} (CDIAL 4345), with the diphthong worn down to \emph{o}. It describes a complexion, as \emph{sā̃vlā} does.
\end{cousinweb}

\begin{grammarlens}[title={What you've built}]
One more word for describing what you lived and felt.
\end{grammarlens}

\subsection*{Your turn}
\begin{itemize}
\raggedright
  \item \emph{Say it:} \emph{gorā}
  \item \emph{Say it:} \emph{gorā}, once more
  \item \emph{Say it:} \emph{Mīrā kā bhāī gorā hai aur Mīrā sā̃vlī hai.} — Mira's brother is fair and Mira is dusky.
  \item \emph{Recall:} say the Hindi for appearance, then the Hindi for dusky, then say \emph{gorā} again
\end{itemize}

\begin{tcolorbox}[breakable,colback=teal!4,colframe=teal!35!black,title={Before you move on}]
What does \hi{गोरा} mean? (\textquotedblleft{}Fair-skinned\textquotedblright{}.) Say it once more.
\end{tcolorbox}
\section[\texorpdfstring{\hi{दुबला} (dublā)}{दुबला (dublā)}]{\hi{दुबला} (dublā) — lean; thin and weak}

\label{lesson:HI-C482-dubla}



\begin{quote}
Before the new one: say the Hindi for dusky, then the Hindi for fair-skinned.
\end{quote}

\subsection*{You'll want to know: \hi{दुबला}}
\textbf{\hi{दुबला}} — \emph{dublā} — \textquotedblleft{}lean; thin and weak\textquotedblright{}.

\begin{quote}
\hi{बीमारी} \hi{के} \hi{बाद} \hi{दादा} \hi{बहुत} \hi{दुबले} \hi{हो} \hi{गए।}
\end{quote}

\begin{quote}
\emph{bīmārī ke bād dādā bahut duble ho gae.}
\end{quote}

\begin{quote}
Grandfather got very thin after his illness.
\end{quote}

\begin{cousinweb}
\textbf{\hi{दुबला}} (\emph{dublā}) is a tadbhava: CDIAL lists Hindi \emph{dublā}, \textquotedblleft{}weak, emaciated\textquotedblright{}, under Sanskrit \emph{durbala}, \textquotedblleft{}weak\textquotedblright{}, from \emph{dur-}, \textquotedblleft{}ill\textquotedblright{}, and \emph{bala}, \textquotedblleft{}strength\textquotedblright{} (CDIAL 6438). \emph{patlā} is slim; \emph{dublā} is thin from weakness.
\end{cousinweb}

\begin{grammarlens}[title={What you've built}]
One more word for describing what you lived and felt.
\end{grammarlens}

\subsection*{Your turn}
\begin{itemize}
\raggedright
  \item \emph{Say it:} \emph{dublā}
  \item \emph{Say it:} \emph{dublā}, once more
  \item \emph{Say it:} \emph{bīmārī ke bād dādā bahut duble ho gae.} — Grandfather got very thin after his illness.
  \item \emph{Recall:} say the Hindi for dusky, then the Hindi for fair-skinned, then say \emph{dublā} again
\end{itemize}

\begin{tcolorbox}[breakable,colback=teal!4,colframe=teal!35!black,title={Before you move on}]
What does \hi{दुबला} mean? (\textquotedblleft{}Lean; thin and weak\textquotedblright{}.) Say it once more.
\end{tcolorbox}
